% ============================================================
% 記入例・模範解答版の追加定義（answers.tex から preamble.tex の後に読み込む）
%   赤字が記入例・模範解答。sessions/NN.tex のマクロに「中身を入れる版」を足す。
% ============================================================
\usepackage{amssymb}

\definecolor{ans}{HTML}{C62828}
\definecolor{anstint}{HTML}{FBE9E7}

% 記入例の文字
\newcommand\ans[1]{{\color{ans}\sffamily #1}}
% 表の中など、小さめの記入例
\newcommand\af[1]{{\color{ans}\sffamily\footnotesize #1}}

% 氏名欄とフッター
\renewcommand*\wsname{\ans{記入例}}
\fancyfoot[L]{\sffamily\scriptsize\color{inkthree}経営情報 2026　AI利用ガイド　復習教材　\textcolor{ans}{記入例}}

% 記入欄（中身入り）　\abox[欄内の小見出し]{最小の高さ}{記入例}
\newcommand\abox[3][]{%
  \begin{tcolorbox}[enhanced,colback=white,colframe=rulegray,boxrule=0.5pt,arc=1.2mm,
      height from=#2 to 260mm,left=2mm,right=2mm,top=1.2mm,bottom=1.5mm,
      fontupper=\sffamily\scriptsize\color{inkthree},before skip=1mm,after skip=1.5mm]
    \if\relax\detokenize{#1}\relax\else #1\par\vspace{0.6mm}\fi
    {\small\color{ans}#3}\end{tcolorbox}}

% 横並びの記入欄（中身入り）　\atwoboxes{見出し1}{見出し2}{高さ}{記入例1}{記入例2}
\newcommand\atwoboxes[5]{%
  \par\noindent
  \begin{minipage}[t]{0.49\linewidth}\abox[#1]{#3}{#4}\end{minipage}\hfill
  \begin{minipage}[t]{0.49\linewidth}\abox[#2]{#3}{#5}\end{minipage}\par}

% 罫線に書き込んだ記入例　\alines{行数}{記入例}
\newcommand\alines[2]{%
  \par\noindent
  \begin{tcolorbox}[enhanced,blankest,height=\dimexpr7.5mm*#1\relax,
      underlay={\foreach\k in {1,...,#1}{\draw[rulegray,line width=0.4pt]
        ([yshift=-\k*7.5mm]frame.north west) -- ([yshift=-\k*7.5mm]frame.north east);}},
      fontupper=\sffamily\color{ans}\fontsize{9.5pt}{7.5mm}\selectfont,
      before skip=0pt,after skip=1.5mm]
    \rule{0pt}{6.4mm}#2
  \end{tcolorbox}}

% 1行の記入欄（中身入り）　\afillin{見出し}{記入例}
%   記入例が長いときは見出しの右で折り返し、最終行の下に罫線を引く。
\newcommand\afillin[2]{%
  \par\noindent\sbox0{\sffamily\small #1\enspace}%
  \hangindent\wd0 \hangafter1 \usebox0\ans{\small #2}\par
  \nointerlineskip\vspace{0.6mm}\noindent\hspace*{\wd0}%
  {\color{rulegray}\rule{\dimexpr\linewidth-\wd0\relax}{0.4pt}}\par\vspace{2.5mm}}

% チェック済みの項目
\newcommand\achk{\makebox[1em][l]{\rlap{□}\hspace{0.12em}\raisebox{0.15ex}{\textcolor{ans}{$\checkmark$}}}}
\newenvironment{achecks}{\begin{itemize}[label=\achk,leftmargin=1.5em,itemsep=0.8mm]\small}{\end{itemize}}

% AI活用記録（中身入り）　\aimemoA{場面}{プロンプト}{扱い}{検証}
\newcommand\aimemoA[4]{%
  \par\needspace{50mm}\vspace{2.5mm}%
  {\sffamily\bfseries\color{phase}■ AI活用記録（この単元の分）}\par\vspace{1mm}
  \begin{wstab}{colspec={Q[l,wd=30mm,m] X[l,m]},row{1-Z}={ht=9mm}}
    使用した場面 & \af{#1} \\
    主なプロンプト & \af{#2} \\
    出力をどう扱ったか\newline{\scriptsize 採用・修正・不採用と理由} & \af{#3} \\
    検証したこと\newline{\scriptsize 何で確かめたか} & \af{#4} \\
  \end{wstab}\par}

% 方眼の記入欄に図を描いた記入例　\agridbox[見出し]{高さ}{TikZ（単位 mm、原点は左下）}
\newcommand\agridbox[3][]{%
  \begin{tcolorbox}[enhanced,colback=white,colframe=rulegray,boxrule=0.5pt,arc=1.2mm,
      height=#2,left=2mm,top=1.2mm,fontupper=\sffamily\scriptsize\color{inkthree},
      underlay={\begin{tcbclipinterior}\draw[step=5mm,black!12,very thin]
        (interior.south west) grid (interior.north east);\end{tcbclipinterior}},
      overlay={\begin{scope}[shift={(interior.south west)},x=1mm,y=1mm,
        every node/.style={font=\sffamily\scriptsize,text=ans,align=center}]#3\end{scope}}]#1\end{tcolorbox}}

% 解説・ポイント
\newenvironment{point}{%
  \par\needspace{20mm}\vspace{3mm}%
  \begin{tcolorbox}[enhanced,colback=anstint,colframe=anstint,boxrule=0pt,arc=1.5mm,
      borderline west={1.2mm}{0pt}{ans},left=4mm,right=3mm,top=1.5mm,bottom=1.5mm,fontupper=\small]
    {\sffamily\bfseries\color{ans}解説・ポイント}\par\vspace{0.5mm}
    \begin{itemize}[label=\textcolor{ans}{・},leftmargin=1.2em,itemsep=0.5mm]}%
  {\end{itemize}\end{tcolorbox}}
